\documentclass{article}

\usepackage[utf8]{inputenc}
\usepackage{amsmath,amssymb}
\usepackage{xurl}
\usepackage[colorlinks=true,linkcolor=blue,citecolor=blue,urlcolor=blue]{hyperref}

% Shared commands for notes in docs/paper-gaps/.

\DeclareUnicodeCharacter{2080}{\ifmmode{}_{0}\else\textsubscript{0}\fi}
\DeclareUnicodeCharacter{2081}{\ifmmode{}_{1}\else\textsubscript{1}\fi}
\DeclareUnicodeCharacter{2082}{\ifmmode{}_{2}\else\textsubscript{2}\fi}
\DeclareUnicodeCharacter{2099}{\ifmmode{}_{n}\else\textsubscript{n}\fi}

\newcommand{\Lean}{\textsc{Lean}}
\ExplSyntaxOn
\NewDocumentCommand{\leanid}{m}{
  \ifmmode
    \textsf{\detokenize{#1}}
  \else
    \group_begin:
    \urlstyle{sf}
    \tl_set:Nn \l_tmpa_tl {#1}
    \tl_replace_all:Nnn \l_tmpa_tl {\_} {_}
    \exp_args:NV \path \l_tmpa_tl
    \group_end:
  \fi
}
\ExplSyntaxOff
\providecommand{\tr}{\operatorname{tr}}
\newcommand{\tnleanIssueBase}{https://github.com/LionSR/TNLean/issues/}
\newcommand{\tnleanPrBase}{https://github.com/LionSR/TNLean/pull/}
\newcommand{\tnleanDocBase}{https://sirui-lu.com/TNLean/docs/}
\newcommand{\ghissue}[1]{\href{\tnleanIssueBase#1}{GitHub issue \##1}}
\newcommand{\ghpr}[1]{\href{\tnleanPrBase#1}{PR \##1}}
\newcommand{\doclink}[1]{\href{\tnleanDocBase#1.html}{\path{#1}}}

% Dirac notation and the identity operator, as in blueprint/src/macros/common.tex.
% \providecommand keeps note-local definitions authoritative.
\usepackage{dsfont}
\providecommand{\ket}[1]{|#1\rangle}
\providecommand{\bra}[1]{\langle #1|}
\providecommand{\braket}[2]{\langle #1 | #2 \rangle}
\providecommand{\ketbra}[2]{|#1\rangle\!\langle #2|}
\providecommand{\Id}{\mathds{1}}
\providecommand{\one}{\mathds{1}}
\providecommand{\C}{\mathbb{C}}
\providecommand{\MN}[1]{M_{#1}(\C)}
\providecommand{\MD}{M_D(\C)}

% External referencing for paper sources.  The build helper writes sanitized
% auxiliary files under build/paper-aux/ and excludes duplicated source labels.
\usepackage{xr}

% Import a generated paper auxiliary file with a caller-supplied label prefix.
% The candidate paths support builds from docs/paper-gaps/, the repository root,
% and build/paper-aux/ itself.
\newcommand{\importpaperaux}[2]{%
  \IfFileExists{../../build/paper-aux/#2.aux}{%
    \externaldocument[#1]{../../build/paper-aux/#2}%
  }{%
    \IfFileExists{build/paper-aux/#2.aux}{%
      \externaldocument[#1]{build/paper-aux/#2}%
    }{%
      \IfFileExists{#2.aux}{%
        \externaldocument[#1]{#2}%
      }{}%
    }%
  }%
}

% General external reference macros
\newcommand{\paperref}[2]{\ref{#1-#2}}
\newcommand{\papereqref}[2]{(\ref{#1-#2})}

% CPSV16 source (1606.00608 / MPDO-22-12-17-2.tex) external document and shortcuts
\importpaperaux{cpsv16-}{1606.00608/MPDO-22-12-17-2-tnlean}
\newcommand{\cpsvref}[1]{\paperref{cpsv16}{#1}}
\newcommand{\cpsveqref}[1]{\papereqref{cpsv16}{#1}}

% Verdict marker: \gapnote{<kind>}{<status>}. Kind is the policy
% classification (clarification, local-correction, scope-restriction,
% unfaithful, false-source, open-gap); status is open, wip (elimination
% underway), resolved, or historical. Machine-read by the site index;
% prints a verdict line.
\newcommand{\gapnote}[2]{\par\noindent\textbf{Verdict:} #1 (#2).\par}


\title{Single-Block Scope for the Schumacher--Werner Commuting-Algebra Proposition}
\author{}
\date{2026-09-01}

\begin{document}
\maketitle
\gapnote{scope-restriction}{open}

\section{The source proposition}

Let $\mathcal H$ be a finite-dimensional complex Hilbert space. Suppose that
$\mathcal A$ and $\mathcal B$ are commuting finite-dimensional subalgebras of
$\mathcal B(\mathcal H)$ with decompositions
\begin{align}
    \mathcal A
     & \cong\bigoplus_{\mu}M_{n(\mu)}(\mathbb C),
    \label{eq:gnvw12_cscom_single_block_scope_left_decomposition} \\
    \mathcal B
     & \cong\bigoplus_{\nu}M_{m(\nu)}(\mathbb C).
    \label{eq:gnvw12_cscom_single_block_scope_right_decomposition}
\end{align}
Proposition~\texttt{Cscom} of Schumacher and Werner states that the product
algebra $\mathcal A\mathcal B$ decomposes into the block-pair products. The
representation on $\mathcal H$ assigns a nonnegative integer multiplicity
$r_{\mu\nu}$ to each pair and satisfies
\begin{align}
    \sum_{\mu,\nu}r_{\mu\nu}n(\mu)m(\nu)
     & =\dim\mathcal H.
    \label{eq:gnvw12_cscom_single_block_scope_multiplicity}
\end{align}
This is Proposition~\texttt{Cscom} and Equation~\texttt{multip} at lines
2119--2138 of Ref.~\cite{SchumacherWerner2004QCA}. For a fixed pair of
blocks, multiplication gives a representation
\begin{align}
    M_{n(\mu)}(\mathbb C)\otimes M_{m(\nu)}(\mathbb C)
     & \longrightarrow\mathcal B(\mathcal H),
    \label{eq:gnvw12_cscom_single_block_scope_block_representation} \\
    A_{\mu}\otimes B_{\nu}
     & \longmapsto A_{\mu}B_{\nu}.
    \label{eq:gnvw12_cscom_single_block_scope_block_multiplication}
\end{align}
The source allows this block-pair representation to vanish because a block
identity need not equal the identity on $\mathcal H$
\cite{SchumacherWerner2004QCA}.

\section{The formal specialization}

The formal result takes one unital block in each factor. Let $I$ be a finite
nonempty set, let
$\mathcal A,\mathcal B\subseteq M_I(\mathbb C)$ be unital star-subalgebras,
and assume
\begin{align}
    \mathcal A
     & \cong M_p(\mathbb C),
    \label{eq:gnvw12_cscom_single_block_scope_left_matrix} \\
    \mathcal B
     & \cong M_q(\mathbb C),
    \label{eq:gnvw12_cscom_single_block_scope_right_matrix}
\end{align}
where $p,q>0$. Since the two subalgebras share the ambient identity, their
multiplication representation is unital and therefore nonzero. Simplicity of
$M_{pq}(\mathbb C)$ makes it injective. Consequently,
\begin{align}
    \mathcal A\vee\mathcal B
     & \cong M_{pq}(\mathbb C),
    \label{eq:gnvw12_cscom_single_block_scope_join} \\
    \dim_{\mathbb C}(\mathcal A\vee\mathcal B)
     & =(pq)^2.
    \label{eq:gnvw12_cscom_single_block_scope_join_dimension}
\end{align}
The formal declarations establish
(\ref{eq:gnvw12_cscom_single_block_scope_join}) and
(\ref{eq:gnvw12_cscom_single_block_scope_join_dimension}).\footnote{Formal
    declarations:
    \leanid{StarSubalgebra.commutingMulMap_injective_of_equiv_matrix},
    \leanid{StarSubalgebra.commutingSupEquivMatrix}, and
    \leanid{StarSubalgebra.finrank_sup_of_equiv_matrix} in
    \doclink{TNLean/Algebra/CommutingStarSubalgebraProduct}.}

This specialization does not formalize the direct sums in
(\ref{eq:gnvw12_cscom_single_block_scope_left_decomposition}) and
(\ref{eq:gnvw12_cscom_single_block_scope_right_decomposition}), the possibility
of zero block-pair representations, or the multiplicity identity
(\ref{eq:gnvw12_cscom_single_block_scope_multiplicity}). Even in the
single-block case, it does not prove the corresponding ambient representation
multiplicity $r$ or the equation
\begin{align}
    r p q
     & =|I|.
    \label{eq:gnvw12_cscom_single_block_scope_single_multiplicity}
\end{align}
Thus the matrix size $pq$ describes the product algebra, not necessarily the
whole ambient matrix algebra.

\section{Why the specialization is useful}

In the adjacent-support-algebra argument, the immediate use of
Proposition~\texttt{Cscom} is that two commuting full matrix factors span a
copy of $M_{pq}(\mathbb C)$; see lines 1270--1308 of
Ref.~\cite{Gross2012IndexTheory}. Showing that this copy fills the adjacent
physical algebra is a separate surjectivity argument. The specialization
(\ref{eq:gnvw12_cscom_single_block_scope_join}) supplies the first statement
without claiming the second.

\section{Elimination plan}

Removing the restriction requires a finite direct-sum form of the
multiplication representation. For every pair $(\mu,\nu)$, restrict
multiplication to
$M_{n(\mu)}(\mathbb C)\otimes M_{m(\nu)}(\mathbb C)$, distinguish the zero and
nonzero representations, and identify each nonzero image with
$M_{n(\mu)m(\nu)}(\mathbb C)$. Decompose $\mathcal H$ into the corresponding
isotypic representation spaces. Their multiplicities then give
$r_{\mu\nu}$, and summing their dimensions proves
(\ref{eq:gnvw12_cscom_single_block_scope_multiplicity}).

Until that construction is formalized, source-labelled statements must mark
the full-matrix, single-block restriction and must not claim the multiplicity
equation.

\bibliographystyle{alpha}
\bibliography{references}

\end{document}